\chaptertitle{第二章\quad 统计图表}

第一章我们把数据整理成了表格。表格虽然准确，但不够直观——要比较各项的大小，还得一行一行看数字。如果把表格变成图画，一眼就能看出谁大谁小、谁高谁低。这一章介绍三种最常用的统计图——条形图、折线图和扇形图——以及如何根据你的需要选择合适的图表。

\sectiontitle{2.1\quad 条形图——比较各类数据的大小}

把第一章中$40$名同学最喜欢的运动项目做成条形图：横轴标出运动项目，纵轴标出人数，每项画一个直条，高度等于对应的频数。

\begin{center}
\begin{tikzpicture}[scale=0.85]
\draw[->,thick] (0,0) -- (5.5,0) node[right]{运动项目};
\draw[->,thick] (0,0) -- (0,4.5) node[above]{人数};
\foreach \y in {0,4,8,12,16} {
  \draw (0.1,\y/4) -- (-0.1,\y/4) node[left]{$\y$};
}
\filldraw[mainblue!80] (0.3,0) rectangle (1.1,4);
\filldraw[mainblue!60] (1.5,0) rectangle (2.3,2.75);
\filldraw[mainblue!40] (2.7,0) rectangle (3.5,1.75);
\filldraw[mainblue!20] (3.9,0) rectangle (4.7,1.5);
\node[mainblue] at (0.7,-0.4){足球};
\node[mainblue] at (1.9,-0.4){篮球};
\node[mainblue] at (3.1,-0.4){乒乓球};
\node[mainblue] at (4.3,-0.4){羽毛球};
\node at (2.5,-1.2){\small 图2-1\quad 最喜欢的运动项目条形图};
\end{tikzpicture}
\end{center}

从这张图一眼就能看出：足球的直条最高（16人），乒乓球和羽毛球的直条较低。这就是条形图的优点——不需要读数字就能比出大小。

\knowbox{
\textbf{条形图}：用高度不同的直条表示各类数据的多少。\\
横轴表示类别（分类轴），纵轴表示数量（数值轴）。\\
直条之间通常留有空隙，表示类别之间是相互独立的。\\
适用于：比较不同类别的数量差异。}

再看一个例子。

\begin{example}
某班期末数学成绩各分数段人数如下，画出条形图描述。\\
90分以上：8人，80~89分：15人，70~79分：10人，60~69分：5人，60分以下：2人。
\end{example}
\begin{solution}
横轴标出五个分数段（从低到高排列），纵轴标人数。直条高度分别为2、5、10、15、8。\\
从图中可以清楚看到"80~89"这一段人最多（15人），"60分以下"人最少（2人）。把分数段按从低到高排列，还能看出"两头少、中间多"的分布特征。
\end{solution}

\begin{center}
\begin{tikzpicture}[scale=0.75]
\draw[->,thick] (0,0) -- (6.5,0) node[right]{分数段};
\draw[->,thick] (0,0) -- (0,4.5) node[above]{人数};
\foreach \y in {0,5,10,15} {
  \draw (0.1,\y/4) -- (-0.1,\y/4) node[left]{$\y$};
}
\foreach \x/\h/\l in {0.3/0.5/<60,1.5/1.25/60~69,2.7/2.5/70~79,3.9/3.75/80~89,5.1/2/$\ge$90} {
  \filldraw[mainblue!60] (\x,0) rectangle (\x+0.8,\h);
  \node at (\x+0.4,-0.4){\l};
}
\node at (3.2,-1.2){\small 图2-2\quad 数学成绩分布条形图};
\end{tikzpicture}
\end{center}

\begin{exercise}
\item 某商店一周内每天的销售额（百元）：周一12，周二8，周三10，周四15，周五20，周六25，周日22。画条形图表示。
\item 条形图中，为什么直条之间要留空隙？如果紧挨着会让人产生什么误解？
\item 图2-2中，如果把分数段的顺序打乱（比如从高到低排列），读图时会不会有什么不同？
\end{exercise}

\sectiontitle{2.2\quad 折线图——观察数据的变化趋势}

条形图适合比较类别之间的大小差异。但如果数据有先后顺序——比如按时间排列——我们更关心的是"变化趋势"：是上升还是下降了？上升得快还是慢？

来看一个例子。某地一年中各月的平均气温如下表：

\begin{center}
\begin{tabular}{c|cccccccccccc}
月份 & 1 & 2 & 3 & 4 & 5 & 6 & 7 & 8 & 9 & 10 & 11 & 12\\\hline
气温($^\circ$C) & 2 & 5 & 10 & 16 & 21 & 26 & 29 & 28 & 23 & 17 & 10 & 4
\end{tabular}
\end{center}

把各月气温标在图上，然后用线段依次连起来：

\begin{center}
\begin{tikzpicture}[scale=0.8]
\draw[->,thick] (0,0) -- (13.5,0) node[right]{月份};
\draw[->,thick] (0,0) -- (0,7.5) node[above]{气温($^\circ$C)};
\foreach \x in {1,...,12} {
  \draw (\x,0.1) -- (\x,-0.1) node[below]{$\x$};
}
\foreach \y in {0,5,10,15,20,25,30} {
  \draw (0.1,\y/4) -- (-0.1,\y/4) node[left]{$\y$};
}
\draw[mainblue,thick,smooth] plot coordinates {
  (1,0.5) (2,1.25) (3,2.5) (4,4) (5,5.25) (6,6.5) (7,7.25) (8,7) (9,5.75) (10,4.25) (11,2.5) (12,1)};
\foreach \x/\y in {1/0.5,2/1.25,3/2.5,4/4,5/5.25,6/6.5,7/7.25,8/7,9/5.75,10/4.25,11/2.5,12/1} {
  \filldraw[mainblue] (\x,\y) circle (2.5pt);
}
\node at (6.5,-1.2){\small 图2-3\quad 某地月平均气温折线图};
\end{tikzpicture}
\end{center}

从折线图可以发现以下信息：
\begin{itemize}
\item 1月到7月气温持续上升，7月最热（29$^\circ$C）；
\item 7月到12月持续下降，1月最冷（2$^\circ$C）；
\item 3月到4月、4月到5月的线段最陡——说明这段时间升温最快；
\item 全年气温呈"先升后降"的单峰形状。
\end{itemize}

\knowbox{
\textbf{折线图}：用点表示各时段的数值，用线段依次连接各点。\\
横轴通常表示时间（或顺序），纵轴表示数值。\\
线段的斜率表示变化的快慢——越陡变化越快。\\
适用于：观察数据随时间或其他顺序的变化趋势。}

\begin{example}
小明上学期$5$次数学测验成绩：第1次72分，第2次78分，第3次85分，第4次82分，第5次90分。画折线图分析他的成绩变化。
\end{example}
\begin{solution}
把次数按顺序放在横轴、分数放在纵轴，标点连线。\\
从折线图看出：整体呈上升趋势（从72→90）。第2次到第3次上升了7分，第4次到第5次上升了8分——这两段进步最快。第3次到第4次略有下降（85→82），但第5次又回到了高点。
\end{solution}

\begin{exercise}
\item 根据图2-3的数据，该地哪个月到哪个月降温最快？你怎么判断的？
\item 小红记录了连续7天的跳绳个数：120、125、130、128、135、140、145。画折线图并分析趋势。
\item 折线图中，相邻两点之间连成的线段，"陡"和"缓"分别代表什么？
\item 条形图和折线图在画法上最大的区别是什么？（提示：直条之间有没有空隙？点之间有没有连线？）
\end{exercise}

\sectiontitle{2.3\quad 扇形图——看各部分占整体的比例}

有时候我们不关心具体数量，只关心"各部分占总体的百分之几"。比如"班上有$40\%$的同学喜欢足球"——这种比例关系用扇形图最合适。

来看一个例子。某班$40$人中不同血型的人数如下：

\begin{center}
\begin{tabular}{c|c|c|c}
血型 & 人数 & 百分比 & 圆心角\\\hline
A型 & 14 & $35\%$ & $126^\circ$\\
B型 & 10 & $25\%$ & $90^\circ$\\
AB型 & 6 & $15\%$ & $54^\circ$\\
O型 & 10 & $25\%$ & $90^\circ$\\
合计 & 40 & $100\%$ & $360^\circ$
\end{tabular}
\end{center}

用扇形图来展示：一个圆代表总体（$100\%$），每个扇形的圆心角等于$360^\circ$乘以对应的百分比。

\begin{center}
\begin{tikzpicture}[scale=1.2]
\draw[thick,fill=mainblue!30] (0,0) -- (0:2) arc (0:126:2) -- cycle;
\draw[thick,fill=exred!30] (0,0) -- (126:2) arc (126:216:2) -- cycle;
\draw[thick,fill=green!30] (0,0) -- (216:2) arc (216:270:2) -- cycle;
\draw[thick,fill=orange!30] (0,0) -- (270:2) arc (270:360:2) -- cycle;
\node at (50:1.3){A型 $35\%$};
\node at (160:1.3){B型 $25\%$};
\node at (240:1.2){AB型 $15\%$};
\node at (315:1.2){O型 $25\%$};
\node at (0,-2.7){\small 图2-4\quad 某班血型分布扇形图};
\end{tikzpicture}
\end{center}

算圆心角的方法：A型占$35\%$，所以圆心角$=360^\circ\times0.35=126^\circ$。B型和O型各占$25\%$，各得$90^\circ$。AB型占$15\%$，得$54^\circ$。所有扇形的圆心角加起来正好是$360^\circ$——这是一个很好的验算方法。

\knowbox{
\textbf{扇形图}：用扇形面积表示各部分占总体的百分比。\\
\textbf{圆心角} $=360^\circ\times$该部分的百分比。\\
所有扇形的圆心角之和等于$360^\circ$（可用于验算）。\\
适用于：显示各部分与整体的比例关系。}

\begin{example}
某家庭月支出：食品$1500$元，房租$1000$元，交通$500$元，教育$800$元，其他$200$元。计算各项占比和圆心角。
\end{example}
\begin{solution}
总支出$=1500+1000+500+800+200=4000$元。\\
食品：$\dfrac{1500}{4000}=37.5\%$，圆心角$=360^\circ\times0.375=135^\circ$。\\
房租：$\dfrac{1000}{4000}=25\%$，$90^\circ$。交通：$\dfrac{500}{4000}=12.5\%$，$45^\circ$。\\
教育：$\dfrac{800}{4000}=20\%$，$72^\circ$。其他：$\dfrac{200}{4000}=5\%$，$18^\circ$。\\
验算：$135+90+45+72+18=360^\circ$。✓
\end{solution}

\begin{exercise}
\item 某班$50$人上学方式：公交20人，步行15人，家长接送10人，骑车5人。求各项百分比和圆心角。
\item 在扇形图中，如果某部分的百分比是$12.5\%$，它的圆心角是多少度？
\item 一个扇形图有4个部分，三个扇形的圆心角分别是$90^\circ$、$120^\circ$、$80^\circ$，第四个扇形的圆心角是多少度？
\end{exercise}

\sectiontitle{2.4\quad 如何选择合适的图表？}

三种图表各有各的用途。有时同一组数据可以用不同的图表来画，但传达的信息不同。

用同一组数据来对比：某班$40$人期末考试成绩分布（90分以上8人，80~89分14人，70~79分10人，60~69分5人，0~59分3人）。

\begin{center}
\begin{tikzpicture}[scale=0.7]
\draw[->,thick] (0,0) -- (6,0) node[right]{分数段};
\draw[->,thick] (0,0) -- (0,4.5) node[above]{人数};
\foreach \y in {0,5,10,14} {
  \draw (0.1,\y/4) -- (-0.1,\y/4) node[left]{$\y$};
}
\foreach \x/\h/\l in {0.3/0.75/0~59,1.2/1.25/60~69,2.1/2.5/70~79,3.0/3.5/80~89,3.9/2/90~100} {
  \filldraw[mainblue!50] (\x,0) rectangle (\x+0.6,\h);
  \node[font=\small] at (\x+0.3,-0.4){\l};
}
\node at (2.5,-1.2){\small (a) 条形图——比较各段人数};
\draw[->,thick] (7.5,0) -- (13,0) node[right]{分数段};
\draw[->,thick] (7.5,0) -- (7.5,4.5) node[above]{人数};
\foreach \y in {0,5,10,14} {
  \draw (7.6,\y/4) -- (7.4,\y/4) node[left]{$\y$};
}
\draw[mainblue,thick] plot coordinates {(7.8,0.75)(8.7,1.25)(9.6,2.5)(10.5,3.5)(11.4,2)};
\foreach \x/\y in {7.8/0.75,8.7/1.25,9.6/2.5,10.5/3.5,11.4/2} {
  \filldraw[mainblue] (\x,\y) circle (2pt);
}
\node[font=\small] at (7.8,-0.4){0~59};
\node[font=\small] at (8.7,-0.4){60~69};
\node[font=\small] at (9.6,-0.4){70~79};
\node[font=\small] at (10.5,-0.4){80~89};
\node[font=\small] at (11.4,-0.4){90~100};
\node at (10.5,-1.2){\small (b) 折线图——观察分布形状};
\end{tikzpicture}
\end{center}

三种图表的对比：

\knowbox{
\textbf{选择指南}：\\
要比较各类的大小 $\rightarrow$ \textbf{条形图}（直条高度比大小）\\
要看变化的趋势 $\rightarrow$ \textbf{折线图}（升降斜率看趋势）\\
要看各部分的占比 $\rightarrow$ \textbf{扇形图}（扇形角度看比例）}

\begin{example}
以下情况各适合用哪种统计图？\\
(a) 比较五个班级的数学平均分；\\
(b) 显示某城市近十年GDP增长趋势；\\
(c) 显示某公司各部门的预算占比。
\end{example}
\begin{solution}
(a) 条形图——比较五个班级谁高谁低。\\
(b) 折线图——看GDP随时间的变化趋势。\\
(c) 扇形图——看各部门预算占总预算的比例。
\end{solution}

\begin{exercise}
\item 以下情况适合用哪种统计图？说明理由。\\
(a) 比较三种品牌手机的市场份额；\\
(b) 展示一周内每天的气温变化；\\
(c) 展示某校各年级的人数分布。
\item 同一组数据能不能既画条形图又画扇形图？它们各自告诉你什么不同的信息？
\item 如果你想强调"某公司今年利润比去年增长了$50\%$"，用哪种图最能突出这个对比效果？画图时有什么需要注意的？
\end{exercise}

\sectiontitle{本章小结}

这一章我们学了三种最常用的统计图。条形图用直条的高度来"比大小"——适合比较不同类别的数量差异。折线图用点的连线来"看趋势"——适合观察数据随时间的变化方向，线段的陡缓反映了变化快慢。扇形图用扇形的角度来"看比例"——适合展示各部分占整体的百分比。选哪种图表取决于你想回答什么问题。三者并不是相互排斥的——同一组数据可以画成不同的图，从不同的角度告诉你不同的信息。

\sectiontitle{课后练习}

\subsection*{A组\quad 基础巩固}

\begin{exercise}
\item 某班同学出生季节统计：春季10人，夏季12人，秋季8人，冬季10人。计算各季度的百分比和扇形圆心角。
\item 下表是某市一周的空气质量指数（AQI）：\\
周一120，周二95，周三105，周四150，周五130，周六80，周日75。\\
画折线图（用文字描述点的位置和趋势），哪一天空气质量最好？哪一天最差？
\item 条形图和折线图最大的画法区别是什么？为什么会有这种区别？
\item 某公司四个季度的销售额（万元）：Q1: 200，Q2: 180，Q3: 220，Q4: 260。\\
选择一种合适的统计图，说明理由。
\item 扇形图中，所有扇形的圆心角之和是多少？为什么？
\item 某班$50$人，近视人数$30$人。用扇形图表示近视与非近视的比例，计算两个扇形的圆心角。
\end{exercise}

\subsection*{B组\quad 挑战提升}

\begin{exercise}
\item 一组数据：A=30，B=20，C=10。画出扇形图时A的圆心角是多少？如果增加一个D=20，A的圆心角会怎么变化？为什么？
\item 某公司全年收入按季度：Q1=100万，Q2=120万，Q3=150万，Q4=180万。\\
(a) 画折线图观察趋势；(b) 计算每个季度的环比增长率；(c) 折线图上哪一段最陡？这和增长率有什么关系？
\item 小明展示零花钱支出扇形图：零食$50\%$，文具$25\%$，交通$15\%$，储蓄$10\%$。\\
(a) 如果每月零花钱$200$元，各项支出是多少？\\
(b) 如果零食降到$30\%$，储蓄会变成多少？扇形图会怎么变？
\item 以下分别用哪种图？\\
(a) 近十年全球平均气温变化；(b) 各国人均GDP排名前10；(c) 某校教师学历构成。
\item 某报纸刊登一张折线图，纵轴从$80$开始而非$0$，显示某公司利润从$80$万涨到$100$万，涨幅看起来非常大。请从$0$开始重画这张图，说明视觉效果有什么不同。这说明了什么统计陷阱？
\end{exercise}

\newpage
